\documentclass[10pt,a4paper]{amsart}
\usepackage[margin=19mm]{geometry}
\usepackage{amsmath,amssymb,mathtools}
\usepackage[scr=rsfs,cal=euler]{mathalfa}
\usepackage{xfrac}
\usepackage[unicode,hidelinks,bookmarks=false]{hyperref}
\newcommand{\ie}{i.e.\ }
\newcommand{\cf}{cf.\ }
\newcommand{\resp}{respectively\ }
\newcommand{\Subsection}{\subsection}
\providecommand{\nobreakdash}{\nobreak\hbox{-}}
\setlength{\emergencystretch}{2em}
\multlinegap0pt
\usepackage{amssymb}
\usepackage{mathtools}
\usepackage[shortlabels]{enumitem}
\newtheorem{theorem}{Theorem}
\newcommand{\E}{\mathbb E}
\newcommand{\F}{\mathbb F}
\newcommand{\R}{\mathbb R}
\DeclareMathOperator{\Tr}{Tr}
\title{\textbf{Euclidean SVP is deterministically NP-hard\\ to approximate within any constant factor}}
\author{Daqing Wan}
\date{Source: arXiv:2608.12664v2 (CC BY 4.0)}
\begin{document}
\maketitle

\noindent\emph{Source and licence.} \textbf{Euclidean SVP is deterministically NP-hard\\ to approximate within any constant factor}. Version arXiv:2608.12664v2 (CC BY 4.0), distributed by arXiv under the Creative Commons Attribution 4.0 International licence, http://creativecommons.org/licenses/by/4.0/, as recorded in the arXiv licence metadata for this exact version and shown on the abstract page https://arxiv.org/abs/2608.12664v2. Full licence notice: \texttt{licenses/arxiv-2608.12664-CC-BY-4.0.txt}.\par\smallskip

\noindent\textit{Editorial scope.} Counted unit: the article's theorem thm:embedding (source0.tex lines 841--849). The article's complete original proof of it is delivered with it (source0.tex lines 850--984, one \texttt{proof} environment of 135 lines). No mathematics is written, repaired or re-derived: the statement and the proof are the article's own lines, in the article's own notation, and no retained line is substituted or re-flowed.  No further source-local context is needed: every cross-reference of every retained line resolves inside the two delivered spans.  Nothing was omitted from any retained span, so the package carries no omission record.  No retained line contains a citation, so the wrapper carries no bibliography and no bibliography entry.  No retained line contains a figure environment, an image-inclusion command or drawing code, so no image is delivered and the package declares no asset dependency.  Editorial formatting only, in addition: an AMS \texttt{amsart} preamble that restates the article's own declarations and macros used by the retained lines, namely the environments \texttt{theorem} on separate counters, together with the standard packages those lines need.  The retained environments print in this document as Theorem 1 in order of appearance, whereas the article prints the same items with the numbering of its own full document.  The complete native source of the article as distributed in the arXiv e-print for this exact version is bundled unchanged as \texttt{sources/207525/source0.tex}.
\begin{theorem}[Enumerated sign embedding]\label{thm:embedding}
Fix $1\le p<\infty$ and an integer $r\ge\max\{2,\lceil p/2\rceil\}$. For every $N\ge2$, a deterministic algorithm constructs $R\in\{-1,1\}^{M\times N}$, where $M=O(N^{2r})$, such that for every $x\in\R^N$,
\begin{equation}\label{eq:embedding}
 \frac{M^{1/p}}{\sqrt3}\|x\|_2
 \le\|Rx\|_p
 \le\sqrt{2r-1}\,M^{1/p}\|x\|_2.
\end{equation}
The algorithm is polynomial-time for fixed $r$, and $R^TR=MI_N$, so $R$ is injective. When $N$ is a power of two, the construction has exactly $M=N^{2r}$ rows. For $1\le p\le2$, the upper constant $\sqrt{2r-1}$ can be replaced by $1$.
\end{theorem}

\begin{proof}
We construct every row explicitly. Averages in the proof are uniform averages over this finite list, counting repeated rows according to their multiplicity. They are a way to evaluate the deterministic construction.

\paragraph{The sign rows.}
Let $Q=2^h$ be the smallest power of two at least $N$, so $h\ge1$ and
\[
 N\le Q<2N.
\]
Choose distinct $t_1,\ldots,t_N\in\F_Q$. We use the absolute field trace
\[
 \Tr(z)=\Tr_{\F_Q/\F_2}(z)
       =\sum_{i=0}^{h-1}z^{2^i}.
\]
This is an $\F_2$-linear map into $\F_2$: additivity follows from the characteristic-two power identities, and $\Tr(z)^2=\Tr(z)$ follows from $z^{2^h}=z$. The displayed polynomial is nonzero and has degree $2^{h-1}<Q$, so it cannot vanish on all of $\F_Q$. The trace is therefore surjective onto $\F_2$. Its kernel has $Q/2$ elements, and each bit has exactly $Q/2$ preimages.

Enumerate all $Q^{2r}$ coefficient tuples of the polynomials
\[
 f(t)=a_0+a_1t+\cdots+a_{2r-1}t^{2r-1},
 \qquad a_i\in\F_Q.
\]
For each coefficient tuple, include the row
\begin{equation}\label{eq:sign-row}
 \bigl((-1)^{\Tr(f(t_1))},\ldots,(-1)^{\Tr(f(t_N))}\bigr)
\end{equation}
in $R$. Thus $M=Q^{2r}<(2N)^{2r}$, which is $O(N^{2r})$ for fixed $r$; if $N$ is a power of two, then $Q=N$ and $M=N^{2r}$. Distinct coefficient tuples may yield the same row; keeping those repetitions makes averaging over rows exactly the same as averaging over all coefficient tuples.

\paragraph{Independence of small sets of coordinates.}
Let $s=(s_1,\ldots,s_N)$ be a uniformly indexed row of $R$. Choose $1\le j\le\min\{2r,N\}$ distinct coordinate positions $i_1,\ldots,i_j$. The evaluation map
\[
 \F_Q^{2r}\longrightarrow\F_Q^j,
 \qquad(a_0,\ldots,a_{2r-1})
 \longmapsto\bigl(f(t_{i_1}),\ldots,f(t_{i_j})\bigr)
\]
is a surjective linear map. Indeed, Lagrange interpolation produces a polynomial of degree at most $j-1<2r$ for any prescribed tuple of evaluations. Its kernel has dimension $2r-j$, so every prescribed tuple has exactly $Q^{2r-j}$ preimages. The $j$ evaluations are consequently independent and uniform on $\F_Q$ under the full enumeration.

Applying the trace coordinatewise gives independent unbiased bits, because each prescribed bit has $Q/2$ preimages. The map $b\mapsto(-1)^b$ is a bijection from $\F_2$ to $\{-1,1\}$, so the corresponding \emph{$j$ signs} in \eqref{eq:sign-row} are independent and unbiased. More explicitly, for every $(\varepsilon_1,\ldots,\varepsilon_j)\in\{-1,1\}^j$,
\[
 \Pr(s_{i_1}=\varepsilon_1,\ldots,s_{i_j}=\varepsilon_j)
 =\frac{Q^{2r-j}(Q/2)^j}{Q^{2r}}=2^{-j}.
\]
This is $2r$-wise independence: every subset of at most $2r$ of the available coordinates has the distribution of independent signs. If $N\ge2r$, the bound on $j$ is $j\le2r$; if $N<2r$, it is $j\le N$, and all $N$ signs are jointly independent. When $N>2r$, full joint independence of all $N$ signs is unnecessary; only the stated independence of small subsets will be used.

Pairwise independence already gives
\begin{equation}\label{eq:embedding-gram}
 R^TR=MI_N.
\end{equation}
To verify this identity entry by entry, write
\[
 (R^TR)_{i\ell}=\sum_{t=1}^M R_{ti}R_{t\ell}.
\]
For $i=\ell$, each summand is $1$, so the sum is $M$. For $i\ne\ell$, pairwise independence says that $(R_{ti},R_{t\ell})$ takes each of the four values $(1,1),(1,-1),(-1,1),(-1,-1)$ on exactly $M/4$ indexed rows. Thus
\[
 (R^TR)_{i\ell}=\frac M4(1-1-1+1)=0.
\]
This proves \eqref{eq:embedding-gram} by exact counting, including any repeated rows.

\paragraph{The lower norm estimate.}
Fix $x\in\R^N$ and set $X=\sum_i s_ix_i$ for the uniformly indexed row $s$ above. As a function of the selected row index, $X(t)=(Rx)_t$ for $1\le t\le M$. Thus, by the definition of a uniform average,
\[
 \E|X|^p=\frac1M\sum_{t=1}^M
       \left|\sum_{i=1}^N R_{ti}x_i\right|^p
       =\frac1M\|Rx\|_p^p.
\]
Consequently,
\begin{equation}\label{eq:finite-moment}
 \|Rx\|_p^p=M\E|X|^p.
\end{equation}
This identity requires no independence; it is simply the sum defining the $\ell_p$ norm written as an average. We now use independence to compare this finite moment with $\|x\|_2$. Since $\E s_i^2=1$ and $\E s_is_j=0$ for $i\ne j$,
\[
 \E X^2=\sum_i x_i^2\E s_i^2
             +2\sum_{i<j}x_ix_j\E s_is_j
          =\sum_i x_i^2=\|x\|_2^2.
\]
Four-wise independence, available since $r\ge2$, gives
\begin{equation}\label{eq:fourth-moment}
 \E X^4=\sum_i x_i^4+6\sum_{i<j}x_i^2x_j^2
 =3\left(\sum_i x_i^2\right)^2-2\sum_i x_i^4
 \le3\|x\|_2^4.
\end{equation}
In these expansions, a term with any sign appearing an odd number of times has average zero. In the fourth moment, the remaining terms either use one index four times, or two indices twice each. For a fixed pair $i<j$, the latter case has $\binom42=6$ arrangements, which explains the coefficient $6$ in \eqref{eq:fourth-moment}. Each expansion uses at most four distinct indices, so this reasoning also applies when $N<4$.

For scalar functions $f,g:\{1,\ldots,M\}\to\R$, H\"older's inequality for the uniform average over row indices states that
\[
 \E|fg|\le(\E|f|^u)^{1/u}(\E|g|^v)^{1/v},
 \qquad u,v>1,\qquad \frac1u+\frac1v=1.
\]
Here the product is pointwise, so $\E|fg|=M^{-1}\sum_{t=1}^M|f(t)g(t)|$. Take $f(t)=|X(t)|^{2/3}$, $g(t)=|X(t)|^{4/3}$ and the conjugate exponents $u=3/2$, $v=3$. Then $fg=|X|^2$, $|f|^{3/2}=|X|$, and $|g|^3=|X|^4$, giving
\[
 \E X^2\le(\E|X|)^{2/3}(\E X^4)^{1/3}.
\]
For $x\ne0$, raising to the power $3/2$ and rearranging gives
\[
 \E|X|\ge\frac{(\E X^2)^{3/2}}{(\E X^4)^{1/2}}
 \ge\frac{\|x\|_2}{\sqrt3}.
\]
For $x=0$, every required inequality holds directly. We also use monotonicity of the means on this probability space. For $0<a<b$, apply H\"older to $|X|^a\cdot1$ with exponents $b/a$ and $b/(b-a)$ to obtain
\[
 \E|X|^a\le(\E|X|^b)^{a/b},
 \qquad (\E|X|^a)^{1/a}\le(\E|X|^b)^{1/b}.
\]
The factor involving the constant function is $1$ because $\E1=1$. Equality of the exponents gives equality of the means. Taking $a=1$, $b=p$ when $p>1$, and using equality when $p=1$, gives
\begin{equation}\label{eq:moment-lower}
 (\E|X|^p)^{1/p}\ge\E|X|\ge\frac{\|x\|_2}{\sqrt3}.
\end{equation}
This establishes a lower bound that includes $p=1$.

\paragraph{The upper norm estimate.}
Since $p\le2r$, monotonicity of the means gives
\[
 (\E|X|^p)^{1/p}\le(\E X^{2r})^{1/(2r)}.
\]
Every term in the expansion of $X^{2r}$ uses at most $2r$ distinct signs, so its average agrees with that for fully independent signs. Only terms in which each index has even multiplicity survive. Each surviving monomial is nonnegative and admits a pairing of its $2r$ positions in which paired positions have the same coordinate index.

A perfect pairing partitions the $2r$ labeled positions into $r$ unordered pairs. Position $1$ has $2r-1$ possible partners. After removing this pair, take the smallest remaining position; it has $2r-3$ possible partners. Continue in this way until the last two positions must be paired. This procedure counts each pairing exactly once, because choosing the smallest unused position fixes the order in which its pairs are formed. The number of perfect pairings is therefore
\[
 (2r-1)(2r-3)\cdots3\cdot1=(2r-1)!!.
\]
For one fixed pairing, summing over an arbitrary index for each pair gives
\[
 \left(\sum_i x_i^2\right)^r=\|x\|_2^{2r}.
\]
Summing this over all pairings counts every surviving term at least once. All terms are nonnegative, so overcounting gives an upper bound:
\[
 \E X^{2r}\le(2r-1)!!\,\|x\|_2^{2r}
 \le(2r-1)^r\|x\|_2^{2r}.
\]
Consequently,
\begin{equation}\label{eq:moment-upper}
 (\E|X|^p)^{1/p}\le\sqrt{2r-1}\,\|x\|_2.
\end{equation}
When $p\le2$, the simpler bound $(\E|X|^p)^{1/p}\le(\E X^2)^{1/2}=\|x\|_2$ improves the upper constant to $1$. Combining \eqref{eq:finite-moment}, \eqref{eq:moment-lower}, and \eqref{eq:moment-upper} proves \eqref{eq:embedding}. Equation~\eqref{eq:embedding-gram} gives injectivity: if $Rx=0$, then $M\|x\|_2^2=x^TR^TRx=0$.

\paragraph{Deterministic construction time.}
The number of rows is $Q^{2r}=O(N^{2r})$, and each has $N$ signs. The field $\F_Q$ can be constructed by finding a monic irreducible polynomial over $\F_2$ of degree $h=\log_2Q$. Even enumerating all monic candidates of this degree and testing divisibility by all monic polynomials of positive degrees at most $h/2$ costs polynomial time in $Q$: there are $Q$ candidates and fewer than $2\sqrt Q$ possible trial divisors. A reducible polynomial has an irreducible factor of degree at most half its degree, so this test is sufficient. An irreducible polynomial of every positive degree exists over $\F_2$. Field arithmetic and polynomial evaluation are then explicit, and the trace uses $h-1$ successive squarings and $h-1$ additions. Since $Q<2N$ and $r$ is fixed, constructing the entire matrix takes deterministic polynomial time in $N$.
\end{proof}

\end{document}
